% Detailed protocol of the application study (moved from the main-text Methods on 2026-10-15, user:
% summary in the Results, details in the SI). Sources as listed in computationalMethodology.tex.
\section{Application protocol: clamping and band gaps}
\label{si:app}

\subsection{Mechanical boundary conditions}
\label{si:brackets-protocol}

\paragraph{Film geometry.}
The $a$ and $b$ axes of the cell are the in-plane directions of a
(001)-oriented film and $c$ is its normal. The stress is
$\sigma = -W/V$, where $W$ is the virial including the kinetic contribution;
positive values are tensile.

\paragraph{Temperature anchoring.}
The absolute transition temperatures of a force field carry its systematic
errors. To place the anneal and the operating point in the same phase and at
the same distance from \Tc{} as in a device, every device temperature $T$ is
mapped onto the model as $T_\mathrm{sim} = T + [\Tc(\mathrm{model}) -
\Tc(\mathrm{expt})]$, using the transition that the anneal-to-operation cooling
path approaches: $\gamma\leftrightarrow\beta$ for \CsSnI{},
$\beta\leftrightarrow\alpha$ for \MASnI{} and the loss of tilt order for
\FASnI{}, each on heating in the production cell. The device temperatures are
the operating temperature (\SI{300}{K}), the damp-heat test (\SI{358}{K}), the
anneal (\SI{373}{K}) and, for \CsSnI{} only, a cubic anneal at \SI{440}{K}. The
shifts are $+49$, \changed{$+27$ and $+4$}~K for \CsSnI{}, \MASnI{} and \FASnI{}, giving
the simulation temperatures of Table~\ref{tab:si-anchored}. A control set at the
unshifted device temperatures (one run per condition) tests the effect of the
anchoring. \changed{The transition temperatures of the hybrids are those of heating runs started from the
model's \SI{50}{K} cell (main text, section \emph{Starting structure of the hybrid heating runs}). The \MASnI{}
shift of $+\SI{27}{K}$ replaces a shift of $-\SI{65}{K}$ obtained from runs started in an oversized cell; the
\MASnI{} runs at 235, 293 and \SI{308}{K} made with that shift are listed in Table~\ref{tab:si-brackets} but not
used in the main text. The \FASnI{} shift was computed from the earlier runs (\SI{242 \pm 4}{K}); the runs from
the \SI{50}{K} cell give \SI{248 \pm 6}{K}, a shift larger by \SI{6}{K}, within the error, so the \FASnI{}
campaigns were kept.}

\begin{table}[htbp]
  \centering
  \caption{Device temperatures and the anchored simulation temperatures
    (K). The shift is $\Tc(\mathrm{model}) - \Tc(\mathrm{expt})$ of the
    anchor transition.}
  \label{tab:si-anchored}
  \begin{tabular}{lcccc}
    \toprule
    & \CsSnI{} & \MASnI{} & \FASnI{} \\
    \midrule
    anchor transition & $\gamma\leftrightarrow\beta$ & $\beta\leftrightarrow\alpha$ & tilt loss \\
    \Tc{} model / expt & 411 / 362 & \changed{302 / 275} & \changed{242$^{b}$ / 238$^{a}$} \\
    shift & $+49$ & \changed{$+27$} & \changed{$+4$} \\
    operation (\SI{300}{K}) & 349 & \changed{327} & \changed{304} \\
    damp heat (\SI{358}{K}) & 407 & \changed{385} & \changed{362} \\
    anneal (\SI{373}{K}) & 422 & \changed{400} & \changed{377} \\
    cubic anneal (\SI{440}{K}) & 489 & -- & -- \\
    \bottomrule
  \end{tabular}
  \begin{flushleft}
    \small $^{a}$ Midpoint of the experimental range \SIrange{225}{250}{K}.\cite{schueller2018crystal}
    \changed{$^{b}$ From the runs started in the oversized cell; \SI{248 \pm 6}{K} from the \SI{50}{K} cell.}
  \end{flushleft}
\end{table}

\paragraph{Three brackets.}
A real film lies between two limits: a free-standing crystal and a crystal
whose in-plane dimensions are fixed by a rigid substrate. We simulate both
limits and one realistic intermediate state, each with three independent runs
at the operating temperature and one at the damp-heat temperature. The free
bracket (ref) applies a constant pressure $P_\mathrm{ext}$ on all three axes.
The stress-controlled bracket (sctl) applies an in-plane tensile stress of
\SI{50}{MPa} ($P_{xx} = P_{yy} = P_\mathrm{ext} - \SI{0.05}{GPa}$) with the
normal free ($P_{zz} = P_\mathrm{ext}$), the residual stress level measured in
perovskite films.\cite{Rolston2018AEM} The rigid clamp with process history
(proc) consists of \SI{50}{ps} of free constant-pressure MD at the anneal
temperature, after which $a$ and $b$ are frozen while $c$ stays free; the cell
is then cooled at \SI{2}{K/ps} to the damp-heat temperature (held
\SI{50}{ps}) and to the operating temperature (held \SI{100}{ps}). The thermal
contraction of the substrate itself is neglected; the residual stress of a
film arises because it is much smaller than that of the
perovskite.\cite{Rolston2018AEM, ref_processing_temp} For \CsSnI{} the same
protocol is also started from the cubic anneal (proccub).
Equilibrium runs start from a heating-ramp frame of the free crystal close to
the target temperature, with the cell set to the five-run mean at that
temperature, and consist of \SI{5}{ps} of constant-volume MD followed by
\SI{150}{ps} at the target condition. The first \SI{50}{ps} of the last stage
are discarded and averages are taken over the remaining frames.

\paragraph{Elastic response.}
Two ladders at the operating temperature give the elastic constants of each
model and the strain dependence of the structure: biaxial strains of $\pm0.25$,
$\pm0.5$, $\pm1$ and $\pm\SI{2}{\percent}$ applied to $a$ and $b$ (frozen,
$c$ free), and hydrostatic pressures of $-0.3$, $-0.15$, $+0.15$, $+0.3$ and
$+\SI{0.6}{GPa}$ relative to $P_\mathrm{ext}$. The biaxial modulus
$Y/(1-\nu)$, Poisson's ratio $\nu$, Young's modulus $Y$ and the bulk modulus
$B$ are fitted to the mean stresses and cells.

\paragraph{Transition temperatures under the clamp.}
Because the model supercools, the shift of a transition under the clamp is
measured on heating. For \CsSnI{}, cells are clamped in-plane at the mean free
cell of the anneal state and heated from the operating temperature at
\SI{2}{K/ps} through both transitions (three runs). The transition temperatures
are determined as for the free ramps, and $\dTc = \Tc(\mathrm{clamped}) -
\Tc(\mathrm{free})$. A $\gamma$-phase cell held under the same clamp at the
operating temperature (\SI{150}{ps}, three runs) provides the clamped $\gamma$
state for the band-gap calculations. Clamped heating runs have so far been made for \CsSnI{} only.

\subsection{Band gaps of the clamped and free ensembles}
\label{si:gaps-protocol}

\paragraph{DFT on production frames.}
Band gaps are computed with the label settings of the main text (\emph{Electronic structure}) on
complete production frames of 2000 or 2304 atoms, so no subcell has to be cut
out of the MD cell. The gap is the lowest unoccupied minus the highest occupied
Kohn--Sham eigenvalue. The supercells sample only the $\Gamma$ point. In the
$\sqrt{2}\times\sqrt{2}\times2$ setting the supercell vectors are
$n_1(1,1,0)$, $n_2(-1,1,0)$ and $n_3(0,0,2)$ in units of \apc{}, so the
pseudo-cubic $R$ point, which carries the band extrema of the cubic phase and
from which those of the tilted phases derive, folds onto $\Gamma$ for every
$n_1$, $n_2$, $n_3$; so does the in-phase tilt vector $M=(\tfrac12,\tfrac12,0)$.
In a pseudo-cubic supercell with an odd number of repeats along an axis, $R$
would instead fold onto the zone boundary, and a $\Gamma$-only gap would be too
large. Each DFT calculation also returns the energy of the frame, which tests
UrsaMLFF in the clamped ensembles.

\paragraph{$\Gamma$-point sampling.}
The folding argument holds for the average structure. A thermal snapshot has
no symmetry beyond the supercell lattice, and disorder could move a band edge
to another point of the supercell zone. Table~\ref{tab:si-gap-kpoints}
collects every DFT calculation of this work that samples more than one
$k$-point. In all 86 static structures and in all 62 MD snapshots of 192 and
240 atoms, both band edges are at $\Gamma$, by at least \SI{213}{meV} for the
valence band and \SI{4}{meV} for the conduction band. The conduction band has
the smaller margin because, without SOC, its minimum at $R$ is threefold
degenerate and includes heavy branches. In the 80-atom fcc-type cell at
\SI{600}{K}, one of two snapshots has its conduction-band minimum off
$\Gamma$, and its $\Gamma$ gap is \SI{11}{meV} too large. Such an error is
bounded by the width of the band-edge bands across the supercell zone, which
falls roughly as the inverse square of the cell length: by a factor of about
five from the 80-atom cell to the hybrid production cells and of about eight
to the \CsSnI{} cells. The $\Gamma$ gap can only overestimate the band-edge
gap, and by similar amounts in clamped and free frames of the same compound,
so the error largely cancels in \dEg{}.

\paragraph{Frame selection.}
Coordinates are never averaged: the gap of the mean structure differs from
the mean gap because averaging removes the tilts that open the gap. Instead, we
select individual frames. For every frame in the averaging window a vector of
structural descriptors is computed: the root-mean-square tilt and the in-phase
and antiphase tilt order about each pseudo-cubic axis, the mean, minimum and
fifth percentile of the Sn--I--Sn angles, the mean, spread and 95th percentile
of the Sn--I bond lengths, the off-centring of Sn in its octahedron, the 95th
percentile of the tilt angles, and for the hybrids the polarisation of the
cation axes, the fractions of cations nearest to $\langle100\rangle$,
$\langle110\rangle$ and $\langle111\rangle$ and the nearest-neighbour
orientational correlation. Within each run the autocorrelation time $\tau$ of
the first principal component of these descriptors defines a spacing
$\Delta = \max(2\tau, \SI{5}{ps})$, and the frames of a bracket are thinned to
at least $\Delta$ apart. The candidates of a bracket are divided into strata,
run $\times$ first or second half of the window, or into clusters when
$k$-means finds a well separated split of the descriptors (silhouette above 0.5
and the smaller cluster above \SI{15}{\percent}), as expected for coexisting
tilt states. Frames are drawn at random within strata with a fixed seed.

\paragraph{Two-level estimator.}
Level~1 isolates the elastic part of the gap change. Six free frames per
material are each computed as they are and after an affine strain to the mean
cell of the stress-controlled bracket, of the rigid clamp, and of the
$\pm\SI{0.3}{GPa}$ hydrostatic runs. Pairing the same configuration cancels
the configurational noise, and the paired differences give the gap change at
fixed internal structure and the deformation potential $\mathrm{d}\Eg/\mathrm{d}\ln V$.
Level~2 compares the ensembles: the mean gap of each bracket is the
stratum-weighted mean $\langle \Eg\rangle = \sum_k w_k m_k$ with the stratum
weights $w_k$ of the candidate pool, and its variance is $\sum_k w_k^2
s_k^2/n_k$; for a stratum with a single frame the spread of the bracket is used.
A pilot of two frames per bracket was planned to set a Neyman
allocation\cite{Neyman1934} of further frames in proportion to $w_k s_k$; its
within-stratum spreads were too uncertain for that, so the three brackets at the
operating temperature were completed to 16 frames each with the frames
distributed over the strata in proportion to $w_k$. Level~1 uses the mean cells
of the first three runs of each bracket, which differ from the eight-run means by
at most \SI{0.08}{\percent} in-plane for the hybrids and \SI{0.16}{\percent}
for the rigid clamp of \CsSnI{}. A ridge regression of the gap on the descriptors was foreseen as a
control variate, $\langle \Eg\rangle \approx \overline{f}_\mathrm{pool} +
\overline{(\Eg - f)}_\mathrm{DFT}$, where $f$ is evaluated on every candidate
frame. On the present frames the frame-averaged descriptors do not predict the
gap (Table~\ref{tab:si-gap-scatter}), so the stratified means are reported
without it. The difference between Level~2 and Level~1 is the part of the gap
change carried by the relaxation of tilts and cation orientations. For each
bracket we report the mean gap shift relative to the free crystal at the same
temperature and the smallest gap of the frames computed, because carriers
collect in the low-gap tail.

{\changedcolor
For \CsSnI{} two further sets complete the estimators. Five Level-1 free
frames are strained to the mean rigid-clamp cell of all eight runs and to the
cell of the clamp set at the cubic anneal; the first gives an elastic part of
\SI{20 \pm 11}{meV}, in agreement with the \SI{19 \pm 7}{meV} obtained at the
three-run cell, and the two cells differ by \SI{-4.3 \pm 0.8}{meV} on the same
frames. One MD frame per run of the free, stress-controlled and rigidly clamped
brackets at the operating temperature (eight each), the three runs of the
$\gamma$ crystal held under the clamp, and two frames each of the cubic-anneal
clamp and of the free crystal at the damp-heat temperature are relaxed with
MACE-MH-1 + D3 (PBEsol labels; D3 as in the production MD) at the fixed mean
cell of the frame's own run over its production window, to a largest force of
\SI{2e-4}{\hartree\per\bohr}. Frames already computed as they are were
preferred, so that 30 of the 31 relaxed frames have the gap of their MD frame
as a partner. Only the atoms move: the relaxed frame keeps the thermal
expansion and the mechanical state of its bracket, and for the rigid clamp the
in-plane cell is the clamp itself. The cell is not relaxed, because the
pressure offset of MACE at the UrsaMLFF cell (\SIrange{0.3}{0.6}{GPa} of
tension in every relaxed frame) would move it away from the ensemble it is
compared with; since every bracket is relaxed in the same way, only differences
between relaxed frames are used. The relaxation keeps the tilt pattern of every
frame, so that the tilt variant of each rigid-clamp run is already present in
its MD frame, and increases the tilt amplitudes by up to about \ang{4}
(Table~\ref{tab:si-gap-relaxed}).\par}
